% Section 2.4 — Quantum computing background
% Copy into your main manuscript .tex file

\subsection{Quantum computing background}
\label{sec:quantum_background}

This section introduces the quantum computing concepts required to
understand the hybrid classical--quantum architecture presented in
Section~\ref{sec:hybrid_network}.

\subsubsection{Qubits and quantum states}

The fundamental unit of quantum information is the \emph{qubit}, a
two-level quantum system whose state is represented by a unit vector in
$\mathbb{C}^2$.  In the computational basis $\{|0\rangle, |1\rangle\}$,
an arbitrary single-qubit state is written as
\begin{equation}
  |\psi\rangle = \alpha\,|0\rangle + \beta\,|1\rangle, \qquad
  |\alpha|^2 + |\beta|^2 = 1,
\end{equation}
where $\alpha, \beta \in \mathbb{C}$ are probability amplitudes.
Geometrically, the state of a single qubit can be visualised as a point
on the Bloch sphere, parameterised by polar and azimuthal angles.  An
$n$-qubit register is described by a statevector in
$\mathbb{C}^{2^n}$,
\begin{equation}
  |\psi\rangle = \sum_{k=0}^{2^n - 1} \alpha_k\,|k\rangle, \qquad
  \sum_{k=0}^{2^n - 1} |\alpha_k|^2 = 1,
\end{equation}
where $|k\rangle$ denotes the $k$-th computational basis state.  The
exponential growth of the Hilbert space dimension with qubit count
underlies both the potential computational advantage of quantum systems
and the cost of their classical simulation.

\subsubsection{Quantum gates}

Quantum computation proceeds by applying \emph{quantum gates}---unitary
transformations that evolve the statevector while preserving its norm.
The gates used in this work fall into two categories.

\paragraph{Single-qubit rotation gates.}
The Pauli rotation gates $R_X$, $R_Y$, and $R_Z$ perform rotations
about the $x$, $y$, and $z$ axes of the Bloch sphere, respectively,
each parameterised by a continuous angle $\theta$:
\begin{equation}
  R_X(\theta) = e^{-i\theta X/2} =
  \begin{pmatrix} \cos\tfrac{\theta}{2} & -i\sin\tfrac{\theta}{2} \\
  -i\sin\tfrac{\theta}{2} & \cos\tfrac{\theta}{2} \end{pmatrix},
\end{equation}
\begin{equation}
  R_Y(\theta) = e^{-i\theta Y/2} =
  \begin{pmatrix} \cos\tfrac{\theta}{2} & -\sin\tfrac{\theta}{2} \\
  \sin\tfrac{\theta}{2} & \cos\tfrac{\theta}{2} \end{pmatrix},
\end{equation}
\begin{equation}
  R_Z(\theta) = e^{-i\theta Z/2} =
  \begin{pmatrix} e^{-i\theta/2} & 0 \\
  0 & e^{i\theta/2} \end{pmatrix},
\end{equation}
where $X$, $Y$, $Z$ denote the Pauli matrices.  The composition
$R_Z(\theta_3)\,R_Y(\theta_2)\,R_X(\theta_1)$ constitutes a general
single-qubit rotation with three degrees of freedom, capable of
reaching any point on the Bloch sphere.

\paragraph{Entangling gates.}
The controlled-NOT (CNOT) gate is a two-qubit gate that flips the
target qubit if and only if the control qubit is in state $|1\rangle$:
\begin{equation}
  \text{CNOT} =
  \begin{pmatrix} 1 & 0 & 0 & 0 \\
                  0 & 1 & 0 & 0 \\
                  0 & 0 & 0 & 1 \\
                  0 & 0 & 1 & 0 \end{pmatrix}.
\end{equation}
The CNOT gate is the primary source of \emph{entanglement} in the
circuit---correlations between qubits that have no classical analogue.
Entanglement enables the quantum state to encode functions that depend
jointly on multiple input features, analogous to non-linear interactions
between neurons in a classical network.

\subsubsection{Parameterised quantum circuits}

A parameterised quantum circuit
(PQC)~\cite{benedetti2019parameterized}, also referred to as a
variational quantum circuit, applies a trainable unitary
$U(\boldsymbol{\theta})$ composed of parameterised gates to the initial
state $|0\rangle^{\otimes n}$, producing a parameter-dependent quantum
state
\begin{equation}
  |\psi(\boldsymbol{\theta})\rangle
  = U(\boldsymbol{\theta})\,|0\rangle^{\otimes n}.
\end{equation}
The output of the circuit is obtained by measuring a Hermitian
observable $\hat{O}$, yielding an expectation value
\begin{equation}
  f(\boldsymbol{\theta})
  = \langle\psi(\boldsymbol{\theta})|\,\hat{O}\,|\psi(\boldsymbol{\theta})\rangle,
\end{equation}
which is a smooth, differentiable function of
$\boldsymbol{\theta}$.  By optimising $\boldsymbol{\theta}$ via
gradient-based methods to minimise a task-specific loss function, the
PQC acts as a trainable function approximator, playing a role analogous
to a layer in a classical neural network.  The expressibility of the PQC---its
ability to represent a wide class of functions---depends on the circuit
depth, qubit connectivity, and gate set~\cite{sim2019expressibility}.

\paragraph{Data encoding.}
To use a PQC for supervised learning or function approximation,
classical input data must be encoded into the quantum state.  Common
strategies include amplitude encoding, basis encoding, and angle
encoding~\cite{schuld2021machine}.  In this work we adopt angle
encoding, where each input feature $x_i$ is mapped to a rotation angle
via $R_Y(\beta_i\, x_i)$, with $\beta_i$ a trainable scaling parameter.
This approach provides a one-to-one mapping between input features and
qubits and is naturally compatible with gradient-based optimisation.

\subsubsection{Circuit architecture}
\label{sec:circuit_architecture}

The PQC employed in this work (Fig.~\ref{fig:pqc_architecture})
operates on $n$ qubits with $L$ variational layers and consists of
three stages:

\begin{enumerate}
  \item \textbf{Data embedding layer.}  Starting from the initial state
    $|0\rangle^{\otimes n}$, each qubit $i$ receives an $R_Y(\beta_i\,
    x_i)$ rotation that encodes the $i$-th component of the classical
    input vector $\mathbf{x} \in \mathbb{R}^n$, scaled by a trainable
    basis parameter $\beta_i$.  This produces the encoded state
    \begin{equation}
      |\psi_0\rangle = \bigotimes_{i=0}^{n-1}
        R_Y(\beta_i\, x_i)\,|0\rangle.
    \end{equation}

  \item \textbf{Variational layers.}  Each of $L$ layers $l = 1,
    \dots, L$ applies two sub-steps.  First, a general single-qubit
    rotation
    $R_Z(\theta^{(3)}_{l,i})\, R_Y(\theta^{(2)}_{l,i})\,
    R_X(\theta^{(1)}_{l,i})$
    is applied to every qubit $i$, introducing $3n$ trainable parameters
    per layer.  Second, a linear CNOT ladder
    $\text{CNOT}(i, i{+}1)$ for $i = 0, \dots, n{-}2$
    entangles neighbouring qubits.  The total number of trainable gate
    parameters is $3\,n\,L$, plus $n$ basis parameters from the
    embedding layer.

  \item \textbf{Measurement layer.}  The Pauli-$Z$ expectation value
    $\langle Z_i \rangle = \langle\psi|\,Z_i\,|\psi\rangle$ is
    evaluated on each qubit.  These $n$ expectation values are averaged
    to produce a single scalar output:
    \begin{equation}
      f(\mathbf{x}, \boldsymbol{\beta}, \boldsymbol{\theta})
      = \frac{1}{n} \sum_{i=0}^{n-1}
        \langle\psi|\, Z_i \,|\psi\rangle
      \;\in [-1,\, 1].
      \label{eq:pqc_output}
    \end{equation}
    The bounded output range $[-1, 1]$ is rescaled by trainable output
    weights and biases in the surrounding classical layers
    (Section~\ref{sec:hybrid_network}).
\end{enumerate}

\subsubsection{Statevector simulation and differentiability}

On near-term quantum hardware, gradients of PQC outputs with respect to
gate parameters are typically computed via the \emph{parameter-shift
rule}~\cite{mitarai2018quantum, schuld2019evaluating}, which evaluates
the circuit at two shifted parameter values per gate.  For an $n$-qubit
circuit with $p$ parameters this requires $2p$ circuit evaluations per
gradient, which becomes expensive for large parameter counts.

In classical simulation, however, the statevector is stored explicitly
as a $2^n$-dimensional complex array, and each gate application is a
matrix--tensor contraction---a standard differentiable operation.  We
exploit this by implementing the entire PQC as native JAX array
operations: rotation gates are $2\!\times\!2$ unitary matrices applied
via \texttt{jnp.tensordot}, and the CNOT gate is realised through index
permutation and conditional flipping of the statevector tensor.

This design embeds the quantum circuit directly within the JAX
computational graph, enabling three key advantages:
\begin{itemize}
  \item \textbf{Automatic differentiation.}  \texttt{jax.grad} and
    \texttt{jax.jvp} propagate gradients through the full hybrid
    pipeline---classical PINN layers, quantum circuit, and loss
    function---in a single backward or forward pass, without
    parameter-shift overhead.
  \item \textbf{JIT compilation.}  The entire training step, including
    the quantum circuit evaluation, is compiled into a single optimised
    XLA program via \texttt{jax.jit}, eliminating Python-level dispatch
    and framework-bridging overhead.
  \item \textbf{Composability.}  The circuit integrates seamlessly with
    JAX transformations such as \texttt{vmap} (batching over subdomains)
    and \texttt{value\_and\_grad} (simultaneous loss and gradient
    computation), which are used extensively in the FBPINNs training
    loop.
\end{itemize}

The trade-off is that statevector simulation scales as
$\mathcal{O}(2^n)$ in memory and $\mathcal{O}(L \cdot n \cdot 2^n)$ in
computation, limiting the approach to moderate qubit counts.  For the
circuit sizes used in this work ($n \leq 8$, $L \leq 6$), this cost is
negligible compared to the classical network and loss evaluation.

To verify that our JAX implementation produces correct results, we
compare parameter gradients $\partial f / \partial\boldsymbol{\theta}$
against PennyLane's \texttt{default.qubit}
backend~\cite{bergholm2018pennylane} across eight circuit configurations
spanning 2--8~qubits and 1--6~layers.  The maximum absolute deviation
across all configurations is below $2 \times 10^{-7}$, consistent with
\texttt{float32} machine precision
(Fig.~\ref{fig:pqc_gradient_agreement}).


% === BibTeX entries (add to your .bib file) ===
%
% @article{benedetti2019parameterized,
%   title={Parameterized quantum circuits as machine learning models},
%   author={Benedetti, Marcello and Lloyd, Erika and Sack, Stefan and Fiorentini, Mattia},
%   journal={Quantum Science and Technology},
%   volume={4},
%   number={4},
%   pages={043001},
%   year={2019}
% }
%
% @article{mitarai2018quantum,
%   title={Quantum circuit learning},
%   author={Mitarai, Kosuke and Negoro, Makoto and Kitagawa, Masahiro and Fujii, Keisuke},
%   journal={Physical Review A},
%   volume={98},
%   number={3},
%   pages={032309},
%   year={2018}
% }
%
% @article{schuld2019evaluating,
%   title={Evaluating analytic gradients on quantum hardware},
%   author={Schuld, Maria and Bergholm, Ville and Gogolin, Christian and Izaac, Josh and Killoran, Nathan},
%   journal={Physical Review A},
%   volume={99},
%   number={3},
%   pages={032331},
%   year={2019}
% }
%
% @article{sim2019expressibility,
%   title={Expressibility and entangling capability of parameterized quantum circuits for hybrid quantum-classical algorithms},
%   author={Sim, Sukin and Johnson, Peter D and Aspuru-Guzik, Al{\'a}n},
%   journal={Advanced Quantum Technologies},
%   volume={2},
%   number={12},
%   pages={1900070},
%   year={2019}
% }
%
% @book{schuld2021machine,
%   title={Machine Learning with Quantum Computers},
%   author={Schuld, Maria and Petruccione, Francesco},
%   year={2021},
%   edition={2nd},
%   publisher={Springer}
% }
